\section{Numerical Experiments}\label{sec:experiments}

This section validates the proposed estimator through a series of
reproducible numerical experiments. The objectives are:

\begin{itemize}
\item To verify \emph{reliability}: $\eta \ge 1$,
\item To assess \emph{sharpness}: $\eta \approx 1$,
\item To demonstrate the \emph{failure of residual-only estimators},
\item To quantify the \emph{dominance of boundary mismatch}.
\end{itemize}

All experiments are performed using the implementation described in
Section~\ref{sec:estimator}.

\subsection{Experimental setup}

\paragraph{PINN architecture.}
We employ fully-connected neural networks with $4$ hidden layers and
$64$ neurons per layer, using $\tanh$ activation functions.

\paragraph{Training.}
Optimization is performed using Adam followed by L-BFGS.
The loss weights are fixed to
\[
w_r = 1, \qquad w_b = 10.
\]

\paragraph{Sampling.}
Interior and boundary collocation points are sampled uniformly.
All experiments use sufficiently dense sampling to ensure convergence
of the training loss.

\paragraph{Estimator computation.}
\begin{itemize}
\item Residuals are evaluated using automatic differentiation.
\item The dual problem is solved using DST on structured domains
and FEM on the L-shaped domain.
\item The lifting problem is solved using Laplace harmonic extension.
\end{itemize}

\paragraph{Error measurement.}
The true error is computed using the exact solution:
\[
|u - \hat{u}|_{H^1(\Omega)} = \|\nabla(u - \hat{u})\|_{L^2(\Omega)}.
\]

\paragraph{Reproducibility and evaluation protocol.}
All experiments are executed using the provided code repository,
ensuring full reproducibility of results. The estimator is computed
exclusively from the trained PINN solution and problem data, without
access to the exact solution.
The exact solution is used solely for evaluation of the true error and
effectivity index, and does not enter the estimator computation.

\subsection{Benchmark results}

Table~\ref{tab:benchmark_results} reports the true error, estimated error,
and decomposition into residual and lifting contributions for all benchmarks.

\begin{table}[h]
\centering
\caption{Estimator performance across benchmarks.}
\label{tab:benchmark_results}
\resizebox{\textwidth}{!}{%
\begin{tabular}{lrrrrr}
\toprule
Benchmark & True energy error & Estimated energy error & Residual & Lifting & $\eta$ \\
\midrule
Poisson (smooth) & 1.1830e-03 & 1.2423e-03 & 3.2535e-05 & 1.2098e-03 & 1.0501 \\
Variable coefficient diffusion & 2.9830e-03 & 3.1848e-03 & 1.0700e-04 & 3.0778e-03 & 1.0676 \\
L-shaped domain singularity & 4.4128e-02 & 5.7611e-02 & 9.2933e-03 & 4.8318e-02 & 1.3056 \\
\bottomrule
\end{tabular}

}
\end{table}

\paragraph{Observations.}

\begin{itemize}
\item \textbf{Reliability.}
In all cases, $\eta \ge 1$, confirming the theoretical bound.

\item \textbf{Sharpness.}
The estimator remains close to the true error, with
$1.05 \le \eta \le 1.31$.

\item \textbf{Boundary dominance.}
The lifting term is consistently the dominant contribution.
For example, in the Poisson problem:
\[
\text{Residual} \ll \text{Lifting}.
\]

\item \textbf{Singular behavior.}
In the L-shaped domain, the estimator remains reliable despite
reduced regularity.
\end{itemize}

\subsection{Ablation study: necessity of boundary lifting}

To quantify the contribution of each component, we compare:

\begin{itemize}
\item Residual-only estimator,
\item Lifting-only estimator,
\item Full estimator.
\end{itemize}

\begin{table}[h]
\centering
\caption{Ablation study of estimator components.}
\label{tab:ablation_results}
\resizebox{\textwidth}{!}{%
\input{tables/ablation_results}
}
\end{table}

\paragraph{Key findings.}

\begin{itemize}

\item \textbf{Residual-only failure.}
The residual-only estimator underestimates the error by orders of magnitude.
For instance, in the smooth Poisson case:
\[
\eta_r \approx 0.0275,
\]
which is completely unreliable.

\item \textbf{Dominance of lifting.}
The lifting-only estimator is already close to the true error,
indicating that boundary mismatch is the primary source of error.

\item \textbf{Full estimator accuracy.}
The full estimator consistently recovers reliability and achieves
$\eta \approx 1$.

\end{itemize}

These results provide strong evidence that \emph{boundary mismatch is the
dominant error mechanism in PINNs with soft boundary enforcement}.

\subsection{Convergence and effectivity}

Figure~\ref{fig:convergence_errors} shows the decay of the true error
and estimator during training.

\begin{figure}[h]
\centering
\includegraphics[width=0.7\textwidth]{convergence_errors.pdf}
\caption{Training loss vs true error and estimator.}
\label{fig:convergence_errors}
\end{figure}

\paragraph{Key insight.}

\begin{itemize}
\item The training loss decreases rapidly,
\item The true error remains significantly larger,
\item The estimator tracks the true error accurately.
\end{itemize}

This confirms that \emph{training loss is not a reliable indicator of
solution accuracy}.

\subsection{Effectivity index behavior}

Figure~\ref{fig:effectivity_benchmarks} shows the effectivity index
across all benchmarks.

\begin{figure}[h]
\centering
\includegraphics[width=0.7\textwidth]{effectivity_benchmarks.pdf}
\caption{Effectivity index across benchmarks.}
\label{fig:effectivity_benchmarks}
\end{figure}

The effectivity index satisfies
\[
1 \le \eta \le 1.31,
\]
demonstrating both reliability and near-optimal sharpness.

\subsection{Discussion}

The experiments support the following conclusions:

\begin{enumerate}[leftmargin=*]

\item \textbf{Residual-only estimators are fundamentally insufficient}
for PINNs due to weak enforcement of boundary conditions.

\item \textbf{Boundary mismatch dominates the error} across all
considered problems.

\item \textbf{The proposed estimator is reliable and sharp}, providing
a practical certification tool for PINN solutions.

\item \textbf{Training loss is misleading}, as it fails to reflect
energy norm accuracy.

\end{enumerate}

These findings highlight a critical gap in current PINN practice and
demonstrate the necessity of rigorous a posteriori error estimation.